\documentclass[10pt,a4paper]{amsart}
\usepackage[margin=19mm]{geometry}
\usepackage{amsmath,amssymb,mathtools}
\usepackage[scr=rsfs,cal=euler]{mathalfa}
\usepackage{xfrac}
\usepackage[unicode,hidelinks,bookmarks=false]{hyperref}
\newcommand{\ie}{i.e.\ }
\newcommand{\cf}{cf.\ }
\newcommand{\resp}{respectively\ }
\newcommand{\Subsection}{\subsection}
\providecommand{\nobreakdash}{\nobreak\hbox{-}}
\setlength{\emergencystretch}{2em}
\multlinegap0pt
\usepackage{tikz}
\usepackage{xcolor}
\newtheorem{lemma}{Lemma}
\title{Glocal Smoothness:\\ Line search and adaptive step sizes can help in theory too!}
\author{Curtis Fox, Aaron Mishkin, Sharan Vaswani,, Mark Schmidt}
\date{Source: arXiv:2506.12648v2 (CC BY 4.0)}
\begin{document}
\maketitle

\noindent\emph{Source and licence.} Glocal Smoothness:\\ Line search and adaptive step sizes can help in theory too!. Version arXiv:2506.12648v2 (CC BY 4.0), distributed by arXiv under the Creative Commons Attribution 4.0 International licence, http://creativecommons.org/licenses/by/4.0/, as recorded in the arXiv licence metadata for this exact version and shown on the abstract page https://arxiv.org/abs/2506.12648v2. Full licence notice: \texttt{licenses/arxiv-2506.12648-CC-BY-4.0.txt}.\par\smallskip

\noindent\textit{Editorial scope.} Counted unit: the article's lemma lem:lambert (source0.tex lines 1001--1024). The article's complete original proof of it is delivered with it (source0.tex lines 1025--1084, one \texttt{proof} environment of 60 lines). No mathematics is written, repaired or re-derived: the statement and the proof are the article's own lines, in the article's own notation, and no retained line is substituted or re-flowed.  No further source-local context is needed: every cross-reference of every retained line resolves inside the two delivered spans.  Nothing was omitted from any retained span, so the package carries no omission record.  No retained line contains a citation, so the wrapper carries no bibliography and no bibliography entry.  No retained line contains a figure environment, an image-inclusion command or drawing code, so no image is delivered and the package declares no asset dependency.  Editorial formatting only, in addition: an AMS \texttt{amsart} preamble that restates the article's own declarations and macros used by the retained lines, namely the environments \texttt{lemma} on separate counters, together with the standard packages those lines need.  The retained environments print in this document as Lemma 1 in order of appearance, whereas the article prints the same items with the numbering of its own full document.  The complete native source of the article as distributed in the arXiv e-print for this exact version is bundled unchanged as \texttt{sources/179884/source0.tex}.
\begin{lemma}
\label{lem:lambert}
Consider the function
\begin{equation}
\label{eq:lambertW}
h(\delta) =
    \frac{1}{4}\log\left(\frac{\Delta_{0}}{\delta}\right) + (\ell_* + \delta)\log\left(\frac{\delta}{\epsilon}\right),
\end{equation}
    with $\Delta_0 \geq \delta \geq \epsilon > 0$ and $\ell_* \geq 0$.
    Define $\xi = 1/4 - \ell_*$. 
    \begin{itemize}
        \item Case 1: if $\xi \leq \epsilon$ the minimizer of this function is $\delta_*=\epsilon$.
        \item Case 2: if $\xi \geq \Delta_0\log(\Delta_0e/\epsilon)$ the minimizer of this function is $\delta_* = \Delta_0$.
        \item Case 3: otherwise the minimizer of this function is 
        \begin{equation}
        \delta_* = \frac{\xi}{\omega},
    \end{equation}
    or equivalently
    \begin{equation}
        \delta_* = \frac{\epsilon}{e}\exp(\omega),
    \end{equation}
        where $\omega$ is the principal branch of the Lambert $W$ function evaluated at $\xi e/\epsilon$ (which is positive since in this case $\xi > \epsilon > 0$).
    \end{itemize}
\end{lemma}

\begin{proof}
We first note that if $\xi \leq 0$ then $\ell_* > 1/4$ and 
 the solution is given by $\delta_*=\epsilon$. Thus in deriving the other situations below we will assume that $\xi > 0$.
%We first note that $h$ is a differentiable convex function of $\delta$ over the domain $\Delta_0 \geq \delta \geq \epsilon$, thus to find a minimizer it is sufficient to check the directional derivatives at the boundaries and for stationary points in the interior.
    Taking the derivative of $h$  with respect to $\delta$ gives
    \begin{align*}
h'(\delta) & = -\frac{1}{4\delta} + \frac{\ell_*}{\delta} + \log\delta + 1 -\log\epsilon\\
& = \log\left(\frac{\delta e}{\epsilon}\right) - \frac{\xi}{\delta},
    \end{align*}
    and we note that both terms are strictly increasing in $\delta$ (for $\xi > 0$). Thus,
%while its second derivative is given by
%\[
%h''(\delta) = \frac 1 \delta + \frac{\xi}{\delta^2}.
%\]
%We note that the derivative is positive for $\xi \geq 0$ and the function is convex over this range. 
% By convexity 
to find a minimizer it is sufficient to check the directional derivatives at the boundaries and for stationary points in the interior.
    \begin{itemize}
    \item Case 1: at the left boundary of $\delta = \epsilon$ we have
    \[
    h'(\epsilon) = 1 - \frac \xi \epsilon.
    \]
We have that $\epsilon$ is a minimizer if $h'(\epsilon) \geq 0$, which happens if $\xi \leq \epsilon$.
\item Case 2: at the right boundary of $\delta = \Delta_0$ we have
\[
h'(\Delta_0) = \log\left(\frac{\Delta_0 e}{\epsilon}\right) - \frac{\xi}{\Delta_0},
\]
and we have $h'(\Delta_0) \leq 0$ if $\xi \geq \Delta_0\log(\Delta_0 e/\epsilon)$
\item Case 3: in this case we have a stationary point in the interior. 
%Note that in this case we can assume $\xi$ is positive since we have $\xi > \epsilon > 0$.
        Equating $h'(\delta)$ with 0 is equivalent to 
\begin{align*}
    \frac{\xi}{\delta} = \log\left(\frac{\delta e}{\epsilon}\right).
\end{align*}
Applying the exponent function to both sides, we have
\begin{align*}
    \exp\left(\frac \xi \delta\right) = \frac{\delta e}{\epsilon}.
\end{align*}
Multiplying both sides by $\xi/\delta$ gives,
\begin{align*}
 \frac{\xi}{\delta}\exp\left(\frac{\xi}{\delta}\right) = \frac{\xi e}{\epsilon}.
\end{align*}
Now in order to find the minimizer of the above equation, we apply the Lambert $W$ function,
\begin{align*}
 W\left(\frac{\xi}{\delta}\exp\left(\frac{\xi}{\delta}\right)\right) = W\left(\frac{\xi e}{\epsilon}\right),
\end{align*}
and applying the defining property ($W(y)\exp(W(y)) = y$ for any $y$) of the Lambert $W$ function on the left gives
\[
\frac{\xi}{\delta} = W\left(\frac{\xi e}{\epsilon}\right),
\]
where we note that $W(\xi e/\epsilon)$ is real and unique since $\xi e/\epsilon > 0$. Solving for $\delta$ gives
\[
\delta = \frac{\xi}{W(\frac{\xi e}{\epsilon})},
\]
or by re-arranging the Lambert equation $W(y)\exp(W(y)) = y$ as $W(y) = y/\exp(W(y))$ we equivalently have
\[
\delta = \frac{\epsilon}{e}\exp\left(W\left(\frac{\xi e}{\epsilon}\right)\right).
\]
\end{itemize}
\end{proof}

\end{document}
